\section{Experiments}

\subsection{Experimental Setup}

\textbf{Dataset.} Our dataset contains 15,000 smell-chemical pairs after augmentation, split 70/15/15 for train/validation/test.

\textbf{Implementation.} We fine-tune bert-base-uncased with learning rate $2 \times 10^{-5}$, batch size 16, and train for 30 epochs with early stopping (patience=5).

\textbf{Baselines.} We compare against:
\begin{itemize}
    \item \textbf{TF-IDF + Logistic Regression}: Traditional ML approach using n-gram features.
    \item \textbf{Rule-based}: Expert-crafted keyword-to-chemical mappings.
    \item \textbf{BERT-base}: Standard BERT without domain adaptation.
\end{itemize}

\subsection{Evaluation Metrics}

We employ standard multi-label and ranking metrics:
\begin{itemize}
    \item Precision@K and Recall@K for K $\in$ \{1, 5, 10\}
    \item Mean Average Precision (MAP)
    \item Normalized Discounted Cumulative Gain (NDCG)
    \item Micro and Macro F1 scores
\end{itemize}

\subsection{Results}

\begin{table}[h]
\centering
\begin{tabular}{lccc}
\toprule
\textbf{Model} & \textbf{P@5} & \textbf{MAP} & \textbf{F1} \\
\midrule
Rule-based & 0.45 & 0.38 & 0.42 \\
TF-IDF + LR & 0.58 & 0.51 & 0.55 \\
BERT-base & 0.71 & 0.64 & 0.68 \\
\textbf{NeoBERT-Smell} & \textbf{0.78} & \textbf{0.72} & \textbf{0.75} \\
\bottomrule
\end{tabular}
\caption{Performance comparison on test set.}
\end{table}

Our model achieves significant improvements over all baselines, demonstrating the effectiveness of transformer-based approaches for this task.
